\section{Historical exploratory measurements}\label{app:historical}
\suppressfloats[t]
These archived measurements lack complete provider-version provenance. Displayed
backend names are experiment aliases. They are retained for traceability and
prompt-condition comparisons, not as verified model rankings or evidence of
a general relation between capability and privacy.

\subsection{Controlled hand-off measurement}\label{app:historical-handoff}
We use 24 authored multi-agent workflows across seven domains, each holding a
sensitive record. A coordinator delegates to an external specialist under a
permissive prompt. Five runs per workflow give 120 trials per configured backend.
Our auditor flags sensitive identifiers or attributes in 78--82\% of trials for
the three open-weight configurations and 46\% for the historical Claude alias
(Table~\ref{tab:wild}). These are detector outputs on controlled tasks, not
independently adjudicated disclosure rates. Cross-purpose hand-offs are purpose
violations by construction; detecting them tests label consistency rather than
an unknown real-world outcome. The central audit records omit raw message text.

\begin{table}[t]\centering\small
\begin{tabular}{llcc}
\toprule
Recipient model & tier & disclosure & raw$\to$center \\
\midrule
Qwen2.5-7B   & small (open) & 82\% (99/120) & \textbf{0} \\
Llama-3.1-8B & small (open) & 81\% (97/120) & \textbf{0} \\
Qwen2.5-14B  & mid (open)   & 78\% (94/120) & \textbf{0} \\
\textbf{Claude (alias)} & API alias & \textbf{46\% (55/120)} & \textbf{0} \\
\bottomrule
\end{tabular}
\caption{Auditor flag rates on 24 authored workflows $\times$ five runs per
configured backend. Model names are archived aliases, not all independently
verified provider versions. Raw message text is not sent to the central auditor;
this does not imply that metadata reveals no information.}
\label{tab:wild}
\end{table}

\subsection{Can the agent simply be told not to?}\label{sec:instr}
The measurement above invites an obvious objection: our coordinator prompt ends
\emph{``include whatever you think is helpful''}, so a reader may reasonably
suspect the leak is our wording rather than the model's behavior. It also leaves
this paper's central claim---that the safeguard has to be structural---as an
assertion. We test both by sweeping the instruction alone, over the same 24
workflows, eleven models---nine open-weight models from five families and two API configurations---and $3{,}168$ trials. Ground truth here is
deliberately \emph{independent of our auditor}: literal identifiers and attribute
terms taken from the source record, matched against what the coordinator actually
wrote, with each pattern checked against its own record so a drifted annotation
fails loudly rather than scoring zero in our favor.

Removing the invitation does not reduce the measured match rate in this
experiment ($96\%\!\to\!98\%$); this contrast does not rule out effects from
other prompt choices. Instructing the agent helps, and the shape of the
help is the finding. Naming the forbidden categories explicitly leaves $82\%$; handing the model the
record's actual policy---owner, permitted recipient, purpose, and the fact that
this recipient is \emph{not} permitted---leaves $68\%$; Claude (alias) under the strongest condition still leaks in $36\%$ of hand-offs, and GPT-4o is the one model that mostly withholds ($14\%$), to which we return below (Table~\ref{tab:instr}).

In the saved runs, instructions reduce identifier matches from $78\%$ to
$35\%$ and sensitive-attribute matches from $89\%$ to $60\%$.
Figure~\ref{fig:redaction} reports the per-configuration contrasts. These
literal-match measurements do not identify the model's reasoning process or
establish a general relation between parameter count and privacy behavior.

\paragraph{Interpretation boundary.} The historical configurations differ in
literal-match and refusal behavior. Missing provider-version evidence prevents
attributing these differences to a specific model release, alignment method or
capability level. The experiment supports only the observed contrasts within
its saved traces and scoring procedure.

Against the separate literal-match scorer, the auditor has recall $0.74$
(716 scorer-positive hand-offs missed across the 3,168 trials) and precision $0.90$.
There are 224 alerts without a literal match. Without independent semantic
adjudication, disagreements may reflect auditor error or limitations of literal
matching. These are proxy agreement statistics, not guaranteed bounds on
semantic precision or recall.

\begin{table}[t]\centering\small
\begin{tabular}{lccccc}
\toprule
Coordinator instruction & leak & 95\% CI & ident. & attrib. & auditor \\
\midrule
permissive & 96\% & [93, 98] & 78\% & 89\% & 77\% \\
invitation removed & 98\% & [96, 99] & 82\% & 89\% & 76\% \\
categories forbidden & 82\% & [78, 86] & 36\% & 80\% & 58\% \\
\textbf{full policy stated} & \textbf{68\%} & [63, 73] & \textbf{35\%} & \textbf{60\%} & 70\% \\
\midrule
\multicolumn{6}{l}{\emph{unverified Claude configuration alone, permissive $\to$ policy}} \\
\quad identifiers & \multicolumn{5}{l}{$44\%\ \to\ \mathbf{1\%}$ \quad (literal-match rate)} \\
\quad attributes & \multicolumn{5}{l}{$40\%\ \to\ 36\%$ \quad (barely moves)} \\
\bottomrule
\end{tabular}
\caption{Sweeping the coordinator's instruction, pooled over eleven models ($792$ trials per row). Scored against ground truth independent of our
auditor. Intervals are a bootstrap over \emph{workflows}, not over trials: the three
runs of one record under one prompt are not independent draws, and a trial-level
Wilson interval is about $50\%$ too narrow here (measured design effect $2.1$--$2.4$). Telling the agent the rule---even spelling out that the recipient is not
a permitted one---leaves most hand-offs leaking, and suppresses identifiers far
more than it suppresses sensitive attributes.}
\label{tab:instr}
\end{table}

% Shared, editable TikZ. Data: codex_pack/figure_data.json.
\begin{widefigure}[t]\centering
\resizebox{\linewidth}{!}{%
\begin{tikzpicture}[x=1mm,y=1mm]
\node[figtitle, anchor=west] at (0,82) {POLICY CHANGES IDENTIFIERS AND ATTRIBUTES DIFFERENTLY};
\node[figtitle, anchor=west] at (42,74) {Identifiers};
\node[figtitle, anchor=west] at (112,74) {Sensitive attributes};
\draw[draw=figamber!50, fill=figamber!3, rounded corners=1mm, line width=.6pt] (0,1) rectangle (178,13);
\draw[figgrid] (43.0,3) -- (43.0,68);
\node[fignote, anchor=north] at (43.0,0) {0};
\draw[figgrid] (57.0,3) -- (57.0,68);
\node[fignote, anchor=north] at (57.0,0) {25};
\draw[figgrid] (71.0,3) -- (71.0,68);
\node[fignote, anchor=north] at (71.0,0) {50};
\draw[figgrid] (85.0,3) -- (85.0,68);
\node[fignote, anchor=north] at (85.0,0) {75};
\draw[figgrid] (99.0,3) -- (99.0,68);
\node[fignote, anchor=north] at (99.0,0) {100};
\draw[figgrid] (113.0,3) -- (113.0,68);
\node[fignote, anchor=north] at (113.0,0) {0};
\draw[figgrid] (127.0,3) -- (127.0,68);
\node[fignote, anchor=north] at (127.0,0) {25};
\draw[figgrid] (141.0,3) -- (141.0,68);
\node[fignote, anchor=north] at (141.0,0) {50};
\draw[figgrid] (155.0,3) -- (155.0,68);
\node[fignote, anchor=north] at (155.0,0) {75};
\draw[figgrid] (169.0,3) -- (169.0,68);
\node[fignote, anchor=north] at (169.0,0) {100};
\node[figlabel, anchor=west] at (1,66) {Phi-3.5 / 3.8B};
\draw[figseries=figrose] (82.76,66) -- (69.32,66);
\filldraw[figdot=figrose] (82.76,66) circle (1.3pt);
\filldraw[fill=figrose,draw=figrose] (69.32,66) circle (1.3pt);
\draw[figseries=figindigo,dashed] (161.16,66) -- (152.76,66);
\filldraw[figdot=figindigo] (161.16,66) circle (1.3pt);
\filldraw[fill=figindigo,draw=figindigo] (152.76,66) circle (1.3pt);
\node[figlabel, anchor=west] at (1,60) {Mistral / 7B};
\draw[figseries=figrose] (92.84,60) -- (83.32,60);
\filldraw[figdot=figrose] (92.84,60) circle (1.3pt);
\filldraw[fill=figrose,draw=figrose] (83.32,60) circle (1.3pt);
\draw[figseries=figindigo,dashed] (166.76,60) -- (161.16,60);
\filldraw[figdot=figindigo] (166.76,60) circle (1.3pt);
\filldraw[fill=figindigo,draw=figindigo] (161.16,60) circle (1.3pt);
\node[figlabel, anchor=west] at (1,54) {Qwen2.5 / 7B};
\draw[figseries=figrose] (94.52,54) -- (82.76,54);
\filldraw[figdot=figrose] (94.52,54) circle (1.3pt);
\filldraw[fill=figrose,draw=figrose] (82.76,54) circle (1.3pt);
\draw[figseries=figindigo,dashed] (166.76,54) -- (163.4,54);
\filldraw[figdot=figindigo] (166.76,54) circle (1.3pt);
\filldraw[fill=figindigo,draw=figindigo] (163.4,54) circle (1.3pt);
\node[figlabel, anchor=west] at (1,48) {Llama3.1 / 8B};
\draw[figseries=figrose] (88.36,48) -- (63.16,48);
\filldraw[figdot=figrose] (88.36,48) circle (1.3pt);
\filldraw[fill=figrose,draw=figrose] (63.16,48) circle (1.3pt);
\draw[figseries=figindigo,dashed] (165.08,48) -- (136.52,48);
\filldraw[figdot=figindigo] (165.08,48) circle (1.3pt);
\filldraw[fill=figindigo,draw=figindigo] (136.52,48) circle (1.3pt);
\node[figlabel, anchor=west] at (1,42) {Gemma2 / 9B};
\draw[figseries=figrose] (94.52,42) -- (66.52,42);
\filldraw[figdot=figrose] (94.52,42) circle (1.3pt);
\filldraw[fill=figrose,draw=figrose] (66.52,42) circle (1.3pt);
\draw[figseries=figindigo,dashed] (165.64,42) -- (143.24,42);
\filldraw[figdot=figindigo] (165.64,42) circle (1.3pt);
\filldraw[fill=figindigo,draw=figindigo] (143.24,42) circle (1.3pt);
\node[figlabel, anchor=west] at (1,36) {Qwen2.5 / 14B};
\draw[figseries=figrose] (92.28,36) -- (64.28,36);
\filldraw[figdot=figrose] (92.28,36) circle (1.3pt);
\filldraw[fill=figrose,draw=figrose] (64.28,36) circle (1.3pt);
\draw[figseries=figindigo,dashed] (165.08,36) -- (153.32,36);
\filldraw[figdot=figindigo] (165.08,36) circle (1.3pt);
\filldraw[fill=figindigo,draw=figindigo] (153.32,36) circle (1.3pt);
\node[figlabel, anchor=west] at (1,30) {Phi-4 / 14B};
\draw[figseries=figrose] (85.56,30) -- (54.76,30);
\filldraw[figdot=figrose] (85.56,30) circle (1.3pt);
\filldraw[fill=figrose,draw=figrose] (54.76,30) circle (1.3pt);
\draw[figseries=figindigo,dashed] (166.76,30) -- (152.76,30);
\filldraw[figdot=figindigo] (166.76,30) circle (1.3pt);
\filldraw[fill=figindigo,draw=figindigo] (152.76,30) circle (1.3pt);
\node[figlabel, anchor=west] at (1,24) {Mistral Small / 24B};
\draw[figseries=figrose] (85.0,24) -- (57.56,24);
\filldraw[figdot=figrose] (85.0,24) circle (1.3pt);
\filldraw[fill=figrose,draw=figrose] (57.56,24) circle (1.3pt);
\draw[figseries=figindigo,dashed] (165.08,24) -- (148.84,24);
\filldraw[figdot=figindigo] (165.08,24) circle (1.3pt);
\filldraw[fill=figindigo,draw=figindigo] (148.84,24) circle (1.3pt);
\node[figlabel, anchor=west] at (1,18) {Qwen2.5 / 32B};
\draw[figseries=figrose] (89.48,18) -- (57.56,18);
\filldraw[figdot=figrose] (89.48,18) circle (1.3pt);
\filldraw[fill=figrose,draw=figrose] (57.56,18) circle (1.3pt);
\draw[figseries=figindigo,dashed] (165.08,18) -- (148.84,18);
\filldraw[figdot=figindigo] (165.08,18) circle (1.3pt);
\filldraw[fill=figindigo,draw=figindigo] (148.84,18) circle (1.3pt);
\node[figlabel, anchor=west] at (1,12) {GPT-4o};
\draw[figseries=figrose] (83.32,12) -- (43.0,12);
\filldraw[figdot=figrose] (83.32,12) circle (1.3pt);
\filldraw[fill=figrose,draw=figrose] (43.0,12) circle (1.3pt);
\node[figlabel,text=figrose, anchor=south west] at (44.0,13.4) {0};
\draw[figseries=figindigo,dashed] (166.76,12) -- (120.84,12);
\filldraw[figdot=figindigo] (166.76,12) circle (1.3pt);
\filldraw[fill=figindigo,draw=figindigo] (120.84,12) circle (1.3pt);
\node[figlabel,text=figindigo, anchor=south west] at (121.84,13.4) {14};
\node[figlabel, anchor=west] at (1,6) {Opus 4.8};
\draw[figseries=figrose] (67.64,6) -- (43.56,6);
\filldraw[figdot=figrose] (67.64,6) circle (1.3pt);
\filldraw[fill=figrose,draw=figrose] (43.56,6) circle (1.3pt);
\node[figlabel,text=figrose, anchor=south west] at (44.56,7.4) {1};
\draw[figseries=figindigo,dashed] (135.4,6) -- (133.16,6);
\filldraw[figdot=figindigo] (135.4,6) circle (1.3pt);
\filldraw[fill=figindigo,draw=figindigo] (133.16,6) circle (1.3pt);
\node[figlabel,text=figindigo, anchor=south west] at (134.16,7.4) {36};
\node[fignote, anchor=west] at (43,-7) {$\circ$ permissive prompt\qquad $\bullet$ full recipient policy\qquad Rates (\%)};
\end{tikzpicture}}
\caption{Matched prompt conditions across eleven models (72 hand-offs per model per condition). Open markers show permissive prompts; filled markers show the record-specific recipient policy. Identifier and attribute rates are plotted separately: the two highlighted API models nearly eliminate identifier matches but retain different attribute rates. Models are ordered by size; the two frontier models are highlighted.}\label{fig:redaction}
\end{widefigure}


\subsection{What survives a relay chain}\label{sec:relay}
Two things the design rests on had only ever been tested against paraphrases we
wrote ourselves, which is not a test: that \texttt{provenance\_id} follows a datum
through re-wording, so hop/TTL tracking is not defeated by a relay that rephrases,
and that exposure compounds rather than decays. We let real agents do the
rephrasing---each forwards the case onward in its own words, four hops deep,
against a hop budget of two---over the same 24 records, three models, 144 chains.

Exposure does not decay. Two of the three models still carry the subject's
identifiers or attributes at hop~4 in \emph{every} chain ($100\%$); the third
declines only from $94\%$ to $79\%$. The intuition that passing a datum through
successive re-writes dilutes it is simply wrong, and the same identifier/attribute
asymmetry appears again from an entirely different mechanism---Mistral-7B's
identifier rate falls across the chain ($90\!\to\!83\%$) while its attribute rate
\emph{rises} ($94\!\to\!96\%$). Meanwhile the tracking holds: provenance survived
all four model paraphrases in $144/144$ chains and the TTL violation fired in
$144/144$, with zero raw content reaching the center. The failure we were looking
for---a datum that keeps leaking while the auditor loses track of it---did not
occur.

\subsection{Is the inference threat real?}
Our detector fires at $k^{*}=2$ hints, but is that where real inferability begins? We
have an independent LLM play the recipient and infer the subject's withheld attribute
from the $K$ benign fragments it received ($K=1,2,3$), scaled to \textbf{16 sensitive
attributes} across health, finance, employment, legal, and behavioral domains
($5$ trials each). In these saved trials, Claude (alias) infers the withheld attribute from two benign
fragments in \textbf{100\% (16/16)} of attributes and from a \emph{single} fragment in
\textbf{81\% (13/16)}. Other configured recipients differ: on the same 16-attribute
set, open-weight recipients infer at $K{=}2$ in 56--69\% (Qwen2.5-7B, Llama-3.1-8B)
and 81\% (Qwen2.5-14B) of attributes, versus 100\% for the Claude alias---inference
differs across configured runs. Missing version provenance prevents a causal
claim about model capability or progress. The default detector stays silent at
$K=1$ even when the recipient returns the target attribute. This is a detection
coverage gap, not evidence of a general no-false-alarm guarantee. A non-firing
recipient is also not proof that the fragments are uninformative. Per-fragment
likelihood ratios in Section~\ref{sec:auditor} permit earlier alerts for stronger
hints, but the ratios require calibration; these observations alone do not
validate them on unseen tasks.
